\section{Part VI: Revealed Preference and the Limits of Self-Report}

\subsection{Verbal Expression vs. Behavioral Sacrifice}

When an interviewee states, ``I would gladly pay ₹50,000 for this program,'' they are offering information about their emotional state and self-image during a friendly interview. They are \textbf{not} providing evidence of a completed ₹50,000 financial transaction.

In a conversation, verbal praise costs zero rupees, zero hours of weekend study, zero reputational risk, and zero trade-offs against family obligations. In an actual commercial transaction, the buyer must surrender scarce cash, sign binding terms, risk public failure, and sacrifice alternative pleasures. Economists term this disparity \textbf{hypothetical bias} [R13]. The practical rule is simple: treat verbal compliments as social politeness; treat the sacrifice of scarce resources as empirical data.

\subsection{Revealed Preference Is Grounded, Not Infallible}

Revealed preference: observing what people actually choose and purchase in the market: is vastly more reliable than ungrounded surveys. Yet even revealed preference can mislead if interpreted naively:
\begin{itemize}[leftmargin=1.2em,itemsep=0.2ex]
  \item A customer buys an inferior course because aggressive marketing misled them about its contents.
  \item An enterprise signs a software contract to exhaust a ``use-it-or-lose-it'' annual budget, yet never deploys the code.
  \item A student enrolls under high-pressure sales tactics, only to demand a refund on day 3.
\end{itemize}
Therefore, analyze the entire behavioral arc: \textbf{acquisition $\rightarrow$ payment $\rightarrow$ active onboarding $\rightarrow$ completion of difficult tasks $\rightarrow$ verified transformation $\rightarrow$ voluntary referral}.

\subsection{The Hierarchy of Behavioral Commitments}

\begin{table}[h]
\centering
\small
\begin{tabularx}{\columnwidth}{>{\RaggedRight\bfseries}p{2.3cm} >{\RaggedRight}X >{\RaggedRight}p{2.1cm}}
\toprule
Commitment Level & What It Demonstrates & What It Does Not Prove \\
\midrule
Social Like / Praise & Passive topic interest & Commercial intent \\
\addlinespace
Consuming Technical Article & Problem awareness & Willingness to pay \\
\addlinespace
Submitting In-Depth Diagnostic & Investment of cognitive labor and time & Financial budget availability \\
\addlinespace
Requesting Scoped Proposal & Formal organizational consideration & Approver sign-off \\
\addlinespace
Seeking Corporate Approval & Active internal champion & Vendor list clearance \\
\addlinespace
Refundable Small Deposit & Tentative interest with low downside & Full-price stamina \\
\addlinespace
Paid Pilot at Market Price & Real transaction under stated conditions & Long-term retention \\
\addlinespace
Full-Price Cohort Participation & Genuine customer acquisition & Learning mastery \\
\addlinespace
Unprompted Referral / Renewal & Sustained realized value and high trust & Universal market fit \\
\bottomrule
\end{tabularx}
\caption{The ladder of observable behavioral commitments.}
\label{tab:ch6_commitment_ladder}
\end{table}

\subsection{Past Behavior, Present Constraints, Future Choice}

Structure field discovery around 3 strictly separated inquiries:
\begin{enumerate}[leftmargin=1.2em,itemsep=0.2ex]
  \item \textbf{Historical behavior:} What did you actually purchase and complete the last time this problem occurred?
  \item \textbf{Present operational constraints:} What are your unchangeable work hours, team commitments, and budget rules today?
  \item \textbf{Hypothetical future choice:} If offered this specific, bounded intervention next Monday, what exact operational criteria would govern your decision?
\end{enumerate}
Never blend these 3 into a single bucket. Historical behavior proves capacity to act; present constraints establish boundaries; hypothetical choices represent proposals awaiting behavioral verification.

\subsection{Money, Time, and Authority as Separate Currencies}

Novice entrepreneurs treat money as the sole barrier to purchase. In practice, customers transact in 3 fundamentally different currencies:
\begin{itemize}[leftmargin=1.2em,itemsep=0.2ex]
  \item \textbf{Cash capital:} Personal savings, bank debt, or departmental corporate budget.
  \item \textbf{Cognitive time and attention:} Evening study hours, focus energy, and missed weekends.
  \item \textbf{Organizational and social reputation:} Standing with one's manager, peer respect, and internal credibility.
\end{itemize}

A senior engineer may have abundant cash to pay ₹1,00,000 without blinking, but literally zero hours of weekly study time. An aspiring student may have 30 hours of free weekly time, but zero rupees. A middle manager may have both money and time, but refuse to risk their corporate reputation sponsoring an unknown startup vendor. Diagnose which currency is binding.

\subsection{Preserving Realistic Purchase Conditions in Tests}

If you wish to test demand for a premium ₹50,000 engineering studio, hosting a free 1-hour Zoom webinar tells you almost nothing about commercial conversion. A free webinar tests whether people enjoy passive listening on a Tuesday evening; it does not test whether they will commit ₹50,000 and 150 hours of intensive programming.

A valid behavioral experiment must preserve the essential friction of the ultimate product:
\begin{itemize}[leftmargin=1.2em,itemsep=0.2ex]
  \item Transparent, non-discounted pricing.
  \item Explicit, non-negotiable prerequisite expectations.
  \item Rigorous weekly work requirements.
  \item Clear terms of withdrawal and assessment.
\end{itemize}
1 single customer who enrolls under realistic, demanding conditions is worth 500 free webinar registrations.

\subsection{Why ``Would You Pay?'' Distorts Pricing Research}

Asking ``What would you pay for this?'' forces the customer into an unnatural guessing game. They anchor on the lowest plausible number, recall consumer app subscriptions, or name a flattering figure to be polite.

Replace speculative pricing questions with historical anchor comparisons:
\begin{itemize}[leftmargin=1.2em,itemsep=0.2ex]
  \item ``What did you pay for the last technical workshop, software license, or consultant you hired?''
  \item ``How was that fee approved, and what alternative was rejected?''
  \item ``If we offered a 2-day architectural evaluation clinic at ₹25,000, what specific deliverable would your team require before approving that expense?''
\end{itemize}

\subsection{Why Funnel Conversion Rates Can Lie}

Suppose a landing page converts 10\% of visitors into cohort applications. That number appears triumphant until a forensic audit reveals the denominator: 80\% of visitors were former university classmates, the application had no technical prerequisite screening, and zero applicants realized the program required 15 hours of live weekend labs.

Always inspect the \textbf{quality of the conversion denominator}:
\begin{itemize}[leftmargin=1.2em,itemsep=0.2ex]
  \item Were prospects recruited through repeatable, cold channels or personal favors?
  \item Did applicants meet mandatory technical prerequisites?
  \item Did applicants understand the financial fee and time burden?
\end{itemize}
A 1\% conversion rate from qualified cold traffic is infinitely superior to a 15\% conversion rate from polite personal contacts.

\subsection{Adoption, Activation, and Realization}

A transaction is merely the beginning of customer understanding, not the end. The customer lifecycle consists of 3 distinct behavioral choices:
\begin{itemize}[leftmargin=1.2em,itemsep=0.2ex]
  \item \textbf{Adoption:} The decision to pay and enroll.
  \item \textbf{Activation:} The first successful completion of a core value-producing action (e.g., passing the first complex evaluation harness).
  \item \textbf{Realization:} The durable acquisition of independent capability that transfers into the customer's live production environment.
\end{itemize}
Many education companies boast high adoption while suffering from catastrophic activation failure. If 60\% of enrolled students never submit assignment 2, the product has failed the customer's job, regardless of whether tuition was non-refundable.

\subsection{Designing Experiments to Disprove Your Favorite Story}

Human nature instinctively searches for evidence that confirms cherished beliefs. A rigorous founder deliberately designs experiments to \textbf{falsify} their core hypotheses.

If you believe that backend engineers are desperate for code reviews on AI evaluation pipelines:
\begin{itemize}[leftmargin=1.2em,itemsep=0.2ex]
  \item Offer a paid, bounded, 1-on-1 diagnostic review of an evaluation pipeline at ₹5,000.
  \item Outreach to 30 qualified engineers who recently discussed evaluation difficulties online.
  \item Explicitly establish the falsification threshold: ``If fewer than 2 qualified engineers agree to pay and schedule a review within 14 days, our hypothesis regarding immediate, acute, self-funded demand is false.''
\end{itemize}
Setting the disproof threshold before running the experiment prevents post-hoc rationalization.
